\documentclass[10pt,a4paper]{amsart}
\usepackage[margin=19mm]{geometry}
\usepackage{amsmath,amssymb,mathtools}
\usepackage[scr=rsfs,cal=euler]{mathalfa}
\usepackage{xfrac}
\usepackage[unicode,hidelinks,bookmarks=false]{hyperref}
\newcommand{\ie}{i.e.\ }
\newcommand{\cf}{cf.\ }
\newcommand{\resp}{respectively\ }
\newcommand{\Subsection}{\subsection}
\providecommand{\nobreakdash}{\nobreak\hbox{-}}
\setlength{\emergencystretch}{2em}
\multlinegap0pt
\usepackage[shortlabels]{enumitem}
\usepackage{amssymb}
\usepackage{mathtools}
\usepackage{bm}
\newtheorem{lemma}{Lemma}
\newcommand{\I}{\mathrm{I}}
\newcommand{\MM}{\mathbf{M}}
\newcommand{\R}{\mathbb{R}}
\newcommand{\VV}{\mathbf{V}}
\renewcommand{\aa}{\bm{a}}
\newcommand{\bbeta}{\bm{\beta}}
\newcommand{\bloss}{\ell}
\newcommand{\diag}[1]{\mathrm{diag}\left(#1\right)}
\renewcommand\geq{\geqslant}
\renewcommand\leq{\leqslant}
\newcommand{\loss}{\mathcal{L}}
\providecommand{\nor}[1]{\bigl\|{#1}\bigr\|}
\newcommand{\rrr}{\bm{r}}
\providecommand{\scal}[2]{\left\langle{#1},{#2}\right\rangle}
\newcommand{\sos}{\bm{s}}
\newcommand{\sosi}[1]{\bm{s}_{#1}}
\newcommand{\uu}{\bm{u}}
\title{Gradient Flow Polarizes Softmax Outputs towards Low-Entropy Solutions}
\author{Aditya Varre, Mark Rofin, Nicolas Flammarion}
\date{Source: arXiv:2603.06248v1 (CC BY 4.0)}
\begin{document}
\maketitle

\noindent\emph{Source and licence.} Gradient Flow Polarizes Softmax Outputs towards Low-Entropy Solutions. Version arXiv:2603.06248v1 (CC BY 4.0), distributed by arXiv under the Creative Commons Attribution 4.0 International licence, http://creativecommons.org/licenses/by/4.0/, as recorded in the arXiv licence metadata for this exact version and shown on the abstract page https://arxiv.org/abs/2603.06248v1. Full licence notice: \texttt{licenses/arxiv-2603.06248-CC-BY-4.0.txt}.\par\smallskip

\noindent\textit{Editorial scope.} Counted unit: the article's lemma lem:convergence-general-loss (source0.tex lines 1295--1300). The article's complete original proof of it is delivered with it (source0.tex lines 1301--1344, one \texttt{proof} environment of 44 lines). No mathematics is written, repaired or re-derived: the statement and the proof are the article's own lines, in the article's own notation, and no retained line is substituted or re-flowed.  Retained uncounted as the source-local context the proof needs: lines 910--913.  Nothing was omitted from any retained span, so the package carries no omission record.  No retained line contains a citation, so the wrapper carries no bibliography and no bibliography entry.  No retained line contains a figure environment, an image-inclusion command or drawing code, so no image is delivered and the package declares no asset dependency.  Editorial formatting only, in addition: an AMS \texttt{amsart} preamble that restates the article's own declarations and macros used by the retained lines, namely the environments \texttt{lemma} on separate counters, together with the standard packages those lines need.  The retained environments print in this document as Lemma 1 in order of appearance, whereas the article prints the same items with the numbering of its own full document.  The complete native source of the article as distributed in the arXiv e-print for this exact version is bundled unchanged as \texttt{sources/195849/source0.tex}.
\begin{align} \label{eq:gradients-general-loss}
    \nabla_{\VV} \loss(\VV, \aa) &=  - \rrr \sos^{\top},  \\
    \nabla_{\aa} \loss(\VV, \aa) &=  - \left( \diag{\sos} - \sos \sos^{\top} \right) \VV^{\top} \rrr.
\end{align}

\begin{lemma}\label{lem:convergence-general-loss}\textbf{Convergence.}
Consider the gradient flow dynamics given by the equations~\eqref{eq:gradients-general-loss} over the loss function $\loss(\VV, \aa)$, 
\begin{align*}
    \frac{d \, \bloss(\bbeta)}{dt} &\leq - \frac{1}{d} \nor{\nabla_{\bbeta} \bloss}^2 
\end{align*} 
\end{lemma}

\begin{proof}
    With $\rrr = - \nabla_{\bbeta} \bloss $, the time derivatives of the parameters are given by, 
    \begin{align*}
        \frac{d\VV}{dt} &=  \rrr \sos^{\top},  \\
        \frac{d\sos}{dt} &=  \left( \diag{\sos} - \sos \sos^{\top} \right)^2 \VV^{\top} \rrr.
    \end{align*}
    Considering the evolution of $\bbeta = \VV \sos$, we have,
    \begin{align*}
        \frac{d \bbeta}{dt} &= \frac{d \VV \sos}{dt} = \frac{d\VV}{dt} \sos + \VV \frac{d\sos}{dt}, \\
        &= \left( \rrr \sos^{\top} \right) \sos + \VV \left( \diag{\sos} - \sos \sos^{\top} \right)^2 \VV^{\top} \rrr, \\
        &= \left[ \nor{\sos}^2 \I +  \VV \left( \diag{\sos} - \sos \sos^{\top} \right) \VV^{\top} \right] \rrr. 
    \end{align*}
    Let us denote the matrix in the brackets as $\MM(\sos, \VV)$, i.e., 
    \begin{align*}
        \MM(\sos, \VV) = \nor{\sos}^2 \I +  \VV \left( \diag{\sos} - \sos \sos^{\top} \right) \VV^{\top}.
    \end{align*}    
    Using this notation, we can write the evolution of $\bbeta$ as follows:
    \begin{align*}
        \frac{d \bbeta}{dt} &=  \MM(\sos, \VV) \rrr = - \MM(\sos, \VV) \nabla_{\bbeta} \bloss. 
    \end{align*}
    Considering the time derivative of the loss function $\bloss(\bbeta)$, we have,
    \begin{align*}
        \frac{d \bloss(\bbeta)}{dt} &= \scal{\nabla_{\bbeta} \bloss}{\frac{d \bbeta}{dt}} = - \scal{\nabla_{\bbeta} \bloss \, \, }{\, \, \MM(\sos, \VV) \, \nabla_{\bbeta} \bloss}. 
    \end{align*}
    Now we will show that the matrix $\MM(\sos, \VV)$ is positive definite. First, 
    \begin{align*}
        \nor{\sos}^2  \geq \sum_{i=1}^{d} \sosi{i}^2 \geq \frac{1}{d} \left( \sum_{i=1}^{d} \sosi{i} \right)^2 = \frac{1}{d}.
    \end{align*}
    The next term is positive semi-definite since for any vector $\uu \in \R^{d}$, we have,
    \begin{align*}
        \uu^{\top} \VV \left( \diag{\sos} - \sos \sos^{\top} \right) \VV^{\top} \uu &= \left( \VV^{\top} \uu \right)^{\top} \left( \diag{\sos} - \sos \sos^{\top} \right) \left( \VV^{\top} \uu \right), \\
        &= \mu^{\top} \left( \diag{\sos} - \sos \sos^{\top} \right) \mu, \quad \text{ where } \mu = \VV^{\top} \uu, \\
        &= \sum_{i=1}^{d} \sosi{i} \mu_i^2 - \left( \sum_{i=1}^{d} \sosi{i} \mu_i \right)^2 \geq 0,
    \end{align*}
    where the last inequality follows from the Cauchy-Schwarz inequality. Hence we have,
    \begin{align*}
        \uu^{\top} \MM(\sos, \VV) \uu \geq \frac{1}{d} \nor{\uu}^2, \quad \forall \, \uu \in \R^{d}.
    \end{align*}
    This shows that the matrix $\MM(\sos, \VV)$ is positive definite. Using the positive definiteness of $\MM(\sos, \VV)$, we have,
    \begin{align*}
        \frac{d \bloss(\bbeta)}{dt} &\leq - \frac{1}{d} \nor{\nabla_{\bbeta} \bloss}^2 \leq 0.
    \end{align*}
    This shows that the loss function $\bloss(\bbeta)$ is non-increasing along the gradient flow dynamics of $\VV$ and $\sos$.
\end{proof}

\end{document}
